The three research questions point in the same direction. The agents differ, but much of the difference lies in the role the agent plays in an issue, and in whether the analysis looks only at issues that were closed. The findings matter to three audiences in different ways.

\textbf{For maintainers.} Issues written by an agent's own bot are, for both agents where a comparison is possible, closed far less often than the agent's other issues (RQ1). For the review agent CodeRabbit, these are follow-ups from its pull-request reviews, often code-level suggestions (RQ3). Projects that enable a review agent should therefore plan for this stream of issues: route them to a person who triages them, group them, or close them automatically after a set time. Otherwise they accumulate as backlog, the kind of bot activity that developers already describe as noise~\cite{wessel2021disturb}. Maintainers should also not read the labels on issues of trigger-based agents such as Sweep as a classification of the work: there, most labels only invoke the agent (RQ1).

\textbf{For agent developers.} An agent that opens its own issues creates work for maintainers, and a large share of these issues stays open. Developers of such agents could make their issues easier to act on, for example by linking each issue to the review comment and the lines of code it comes from, by stating a priority, or by closing an issue once the underlying code has changed. Agents could also mark clearly how they are involved in an issue. That would reduce the ambiguity that currently forces datasets to attribute issues with heuristics~\cite{ghaleb2026codage,ghaleb2026fingerprinting}.

\textbf{For researchers.} ``AI coding agent'' is not one kind of actor. In the hand-checked sample, the agent had written fewer than a third of the issues correctly attributed to it, and in half it was only mentioned. Four consequences follow for studies that compare agents. First, they should record and report the role of the agent in each artefact, because pooling issues an agent wrote with issues that merely mention it can produce differences that belong to the role, not the agent (RQ1). They should also validate the attribution by hand on a sample; in this study, most wrong attributions came from text signals. Second, they should compare speed with survival analysis: in this study, the ranking of agents by closed issues barely agreed with the censoring-aware ranking, and for one agent the closed issues even reversed a comparison (RQ2). Third, they should remove agent-workflow labels before treating labels as data, adding to known problems with label quality~\cite{herzig2013bug}. Fourth, they should test whether their topics are reproducible, check the topic labels on sample issues, and confirm the results with a second method (RQ3): on this corpus, only two of five fine topics passed both the stability test and the label check. Single LDA runs were not reproducible on this corpus, which echoes earlier warnings~\cite{agrawal2018wrong}, and two reasonable topic models divided the same issues differently while agreeing on the agent-level conclusions.
